\section{Sensitivity analysis}
\label{sec:sensitivity}

\begin{keynote}
\noindent\textbf{This section specifies the design and reports its results.}

The design below was published at methodology version 2026.07.30, before the
calculation engine existed, precisely so that it could be criticised before
execution and so that the results could not be shaped after the fact to suit
the outcome. The analysis was executed in Phase 2 against the twenty
hand-verified golden scenarios; \Cref{sec:sensitivity-results} reports the
results at version \methver{}, produced by
\texttt{tools/sensitivity\_analysis.py} and regenerated whenever any
aggregation weight changes.
\end{keynote}

\subsection{Why sensitivity analysis is obligatory here}

\textcite{nardo2008} treats uncertainty and sensitivity analysis as integral to
a credible composite indicator rather than as an optional addition. The
argument applies with particular force to this methodology.

The aggregation weights of \Cref{eq:aggregation} are the tool's most
consequential unsourced choice. Every other quantitative element --- the
typology envelopes, the canopy threshold, the energy sensitivity, the LST
normalisation --- traces to a published source that a reviewer can consult and
dispute on its own terms. The weights trace to a judgment. That judgment
determines the headline number, and therefore determines which projects a city
funds.

Two responses to this are available. One is to argue that the weights are
reasonable and hope the reader agrees. The other is to measure how much the
output actually depends on them and publish the answer, including if the answer
is inconvenient. Only the second is checkable, and it is the one adopted.

There is a further reason specific to this tool. If rank ordering turns out to
be robust to weight perturbation, then the disclosed equity-forward stance of
\Cref{sec:effective-weights} --- while remaining a genuine value choice ---
changes fewer decisions than its prominence suggests, and deployments can adopt
the tool without first having to resolve a philosophical dispute. If ordering
turns out to be fragile, users need to know that a different but equally
defensible weight set would have produced a different priority list, and the
tool should say so at the point of use. Both outcomes are informative; only
running the analysis distinguishes them.

\subsection{Design of the analysis}

Each aggregation weight in \Cref{eq:aggregation} is varied by
$\pm\,\SI{25}{\percent}$ of its nominal value, with the remaining weights
renormalised to preserve a unit sum, and the full golden-scenario set is
re-scored under each perturbation. Formally, for the perturbation of weight
$w_k$ by a factor $\lambda \in \{0.75,\, 1.25\}$:

\begin{equation}
\label{eq:perturbation}
w_k' = \frac{\lambda w_k}{\lambda w_k + \sum_{i \neq k} w_i},
\qquad
w_j' = \frac{w_j}{\lambda w_k + \sum_{i \neq k} w_i}\ \ \text{for } j \neq k .
\end{equation}

The golden-scenario set is the same collection of hand-verified
input/output cases that serves as the engine's regression suite
(\Cref{sec:testing}), so the scenarios exercised by the sensitivity
analysis are the same ones whose correctness is independently asserted.

\subsection{What is reported}
\label{sec:sa-reported}

Four quantities are reported for each perturbation, chosen because they
correspond to four distinct ways the tool could mislead a user.

\begin{enumerate}
  \item \textbf{Rank stability} --- the proportion of scenario pairs whose
  relative ordering is preserved under the perturbation. This is the primary
  measure. Prioritisation is the tool's actual function, so ordering matters
  more than absolute values: a perturbation that shifts every score by five
  points but reorders nothing has not changed a single decision.

  \item \textbf{Score displacement} --- the distribution of the absolute change
  in the Opportunity Score across the scenario set, reported as a distribution
  rather than a mean, because a small average displacement concealing a few
  large ones is a materially different result from a uniformly small one.

  \item \textbf{Category migration} --- how often a scenario crosses a band
  boundary in \Cref{tab:hpi-bands}, for instance from \emph{Strong
  opportunity} to \emph{Moderate opportunity}. Categories drive the
  recommendation text, so a category change is a change in what the tool
  \emph{says}, not merely in what it scores. A score displacement of two points
  matters greatly at 80.5 and not at all at 45.

  \item \textbf{Influence ranking} --- which weights the output is most
  sensitive to, ordered. This tells a deployment considering recalibration
  which weight to argue about first, and it identifies which of the tool's
  value judgments are actually load-bearing.
\end{enumerate}

\subsection{Results at version \methver{}}
\label{sec:sensitivity-results}

The full scenario set (twenty scenarios, 190 scenario pairs) was re-scored
under each of the twelve perturbations --- six weights, each at
$\lambda \in \{0.75, 1.25\}$ of \Cref{eq:perturbation}.

\begin{table}[H]
\centering
\caption[Sensitivity analysis results]{Sensitivity of the NbS Cooling
Opportunity Score to \SI{\pm 25}{\percent} perturbation of each aggregation
weight, across the twenty golden scenarios. Rank stability is the proportion of
the 190 scenario pairs whose ordering is preserved; displacement is the
absolute change in the Opportunity Score (0--100 scale); migrations counts
scenarios whose interpretation band changed.}
\label{tab:sensitivity-results}
\small
\begin{tabular}{@{}llrrrr@{}}
\toprule
\textbf{Weight} & \textbf{Shift} & \textbf{Rank stab.} &
\textbf{Mean disp.} & \textbf{Max disp.} & \textbf{Migrations} \\
\midrule
Heat Priority Index & $-25\%$ & 0.9737 & 0.56 & 2.00 & 1 \\
                    & $+25\%$ & 0.9895 & 0.49 & 1.77 & 0 \\
Cooling Potential   & $-25\%$ & 0.9895 & 0.63 & 2.01 & 1 \\
                    & $+25\%$ & 0.9842 & 0.55 & 1.77 & 1 \\
NbS Suitability     & $-25\%$ & 0.9895 & 0.47 & 1.49 & 0 \\
                    & $+25\%$ & 0.9842 & 0.44 & 1.38 & 0 \\
Vulnerability       & $-25\%$ & 0.9895 & 0.52 & 1.75 & 0 \\
                    & $+25\%$ & 0.9895 & 0.47 & 1.62 & 0 \\
Co-benefits         & $-25\%$ & 0.9947 & 0.32 & 0.76 & 0 \\
                    & $+25\%$ & 0.9895 & 0.30 & 0.72 & 0 \\
Cost Feasibility    & $-25\%$ & 0.9895 & 0.17 & 1.12 & 0 \\
                    & $+25\%$ & 0.9842 & 0.16 & 1.06 & 0 \\
\bottomrule
\end{tabular}
\end{table}

The four committed quantities, in the order of \Cref{sec:sa-reported}:

\begin{enumerate}
  \item \textbf{Rank stability.} Worst case \num{0.9737}, under the $-25\%$
  Heat Priority Index perturbation: five of 190 pairs reorder. Nine of the
  twelve perturbations preserve at least \SI{98.4}{\percent} of pair
  orderings. Every reordering involves a pair whose baseline scores differ by
  less than about two points --- pairs the methodology itself would describe as
  materially equivalent.

  \item \textbf{Score displacement.} Pooled over all 240
  scenario-perturbation combinations: mean \num{0.42} points, median
  \num{0.32}, 75th percentile \num{0.59}, maximum \num{2.01} on the 0--100
  scale.

  \item \textbf{Category migration.} Three of 240 combinations cross a band
  boundary, each from a baseline score within about \num{1.3} points of the
  boundary (59.58 $\rightarrow$ 60.3 and 60.39 across the Moderate--Strong
  line; 80.92 $\rightarrow$ 79.63 across the Strong--High-priority line). No
  scenario moves by more than one band, and no scenario distant from a boundary
  migrates.

  \item \textbf{Influence ranking.} By mean absolute displacement: Cooling
  Potential (0.59), Heat Priority Index (0.53), Vulnerability (0.50), NbS
  Suitability (0.45), Co-benefits (0.31), Cost Feasibility (0.16). Cost
  feasibility ranks last partly by construction: it is excluded and its weight
  redistributed in the fifteen of twenty scenarios that supply no cost data, so
  the weight binds only where economic evidence exists.
\end{enumerate}

\paragraph{Interpretation.} Under \SI{\pm 25}{\percent} perturbation of any
single aggregation weight, the tool's priority ordering is substantially
stable: at least \SI{97.4}{\percent} of pairwise orderings survive every
perturbation, scores move by well under half a point on average, and category
changes occur only for sites already sitting on a band boundary. The disclosed
equity-forward stance of \Cref{sec:effective-weights} therefore changes fewer
decisions than its prominence suggests: a deployment that rebalances the
weights within this range should expect a materially similar priority list
unless two options were near-tied to begin with --- and near-ties are exactly
the cases the methodology already tells users not to resolve on the headline
score alone.

Two qualifications are owed. The analysis perturbs one weight at a time, so
these figures are not bounds for simultaneous re-weightings. And the golden
scenarios are designed to span the methodology's behaviour --- climates,
typologies, data completeness, cost availability --- not to be a probability
sample of real assessments; the stability figures are properties of this set.
Both qualifications argue for reading near-boundary results as ranges rather
than verdicts, which the confidence mechanism already enforces at the point of
use.

\subsection{Publication and maintenance}

The results above are also committed to the repository in machine-generated
form (\texttt{docs/methodology/SENSITIVITY-ANALYSIS.md}) and in the
Methodology Report, and must be regenerated whenever any weight changes. Because a
weight change requires a methodology version bump
(\Cref{sec:governance}), and because the version stamp is enforced across all
configuration files by continuous integration, a weight change that does not
carry a regenerated sensitivity analysis is visible in review rather than
silent.

The threshold at which a result would prompt a methodology change is stated in
advance to avoid its being negotiated afterwards: if rank stability under any
single $\pm\,\SI{25}{\percent}$ weight perturbation falls below a level at
which the tool's prioritisation could be considered reliable, the appropriate
response is to report that limitation prominently at the point of use ---
telling the user that their priority ordering is weight-sensitive --- rather
than to adjust the weights until the number improves.
